Este capítulo retoma os objetivos, os resultados técnicos apresentados nesta revisão para o TCC~II e as limitações que ainda condicionam a avaliação da plataforma.

\section{Retomada dos Objetivos}

O objetivo geral foi projetar, implementar e integrar uma IDE web ao núcleo de tradução e execução configurável. Os fluxos apresentados no Capítulo~\ref{cap:metodologia} documentam o assistente de personalização, a edição de código, a execução e depuração, a gestão de linguagens e os recursos educacionais de turmas e exercícios.

A especificação dos elementos configuráveis está consolidada na Tabela~\ref{tab:elementos-configuraveis} e na descrição das variantes de análise. A implementação do assistente e dos recursos de edição e execução é apresentada por trechos de código e capturas de tela. A integração utiliza o pacote \texttt{compiler} no navegador e, para a validação de submissões, no servidor Next.js, conforme a Figura~\ref{fig:arquitetura-plataforma}.

A disponibilização técnica permite executar os fluxos locais no navegador. Nesta revisão, 120 testes selecionados pela configuração de integração da IDE foram aprovados, e a Seção~\ref{sec:acessibilidade-responsividade} registra uma verificação delimitada da interface. Os resultados não permitem considerar concluídos os requisitos de acessibilidade, contraste e internacionalização integral, nem demonstram a disponibilidade de uma implantação pública específica.

O objetivo de propor um protocolo de avaliação foi atendido pela Seção~\ref{sec:protocolo-avaliacao}, que relaciona tarefas, decisões de interface, medidas e instrumentos de coleta. A proposta ainda requer definição da amostra, revisão dos instrumentos e encaminhamento institucional antes de sua aplicação. Não se trata de uma avaliação já realizada.

\section{Limitações}

A principal limitação é a ausência de dados de estudantes. O artefato permite investigar a personalização da linguagem, mas ainda não responde empiricamente em que medida essa personalização reduz a barreira sintática. Não se apresentam resultados de aprendizagem, redução de carga cognitiva ou aumento de engajamento.

A verificação de acessibilidade abrange um recorte do assistente, em um navegador, com testes automatizados e uma transição por teclado. Os problemas encontrados exigem ajustes na aplicação; faltam a inspeção dos processos completos e a verificação com tecnologias assistivas. O uso de bibliotecas de componentes não sustenta uma presunção de conformidade.

A configuração de testes utilizada não inclui todos os arquivos existentes, e o fluxo de implantação versionado não demonstra execução automática do Vitest. Também não foram conduzidas medições de desempenho nem auditoria de segurança. A configuração do token em \textit{cookie} acessível por JavaScript e a ausência de isolamento de processo na avaliação de programas são aspectos que exigem análise técnica específica antes de ampliar o cenário de uso.

A linguagem oferece construções imperativas e variantes de tipagem e de arranjos implementadas no núcleo. Não permite acrescentar gramáticas arbitrárias, orientação a objetos ou modularização de um programa em múltiplos arquivos de código-fonte. O sistema de arquivos virtual da IDE organiza documentos, sem representar, por si só, um mecanismo de módulos da linguagem.

\section{Encaminhamentos}

Os próximos passos técnicos são corrigir os problemas de acessibilidade documentados, verificar os fluxos completos de teclado e leitor de tela, revisar os textos ainda não internacionalizados e integrar a execução dos testes ao processo de incorporação de mudanças. As capturas históricas da implementação devem ser atualizadas conforme a interface se estabilize para a apresentação final.

No plano metodológico, a continuidade consiste em revisar e aplicar o protocolo proposto, depois de definir participantes, instrumentos e procedimentos institucionais. Só então será possível discutir resultados de uso e investigar a contribuição pedagógica da personalização. O compartilhamento de linguagens, já implementado no catálogo comunitário e na vinculação a exercícios e listas, poderá ser estudado como recurso didático em turmas; a exportação para uso fora da plataforma permanece outro possível desdobramento.
